\begin{definition}\label{def:reduced_collatz_step}\lean{CollatzMapBasics.reduced_collatz_step}\leanok
The reduced Collatz map $R: \mathbb{N}\to \mathbb{N}$ is defined by
\begin{equation}
    R(n) = \frac{3n+1}{2^{v_2(3n+1)}}
\end{equation}
where $v_2(m)$ is the dyadic valuation of $m$. In short $R(n)=$ odd part of $3n+1$.
\end{definition}

\begin{definition}\label{def:R_iter}\lean{CollatzMapBasics.R_iter}\leanok
For any $k \in \mathbb{N}$, the $k$-fold iteration of the reduced Collatz map
$R^k: \mathbb{N} \to \mathbb{N}$ is defined recursively by
\begin{equation}
    R^0(n) = n
\end{equation}
and
\begin{equation}
    R^{k+1}(n) = R^k(R(n)).
\end{equation}
\end{definition}

\begin{lemma}\label{lemma:R_zero}\lean{CollatzMapBasics.R_zero}\leanok
$R(0) = 1$.
\end{lemma}
\begin{proof}\leanok
By definition, $R(0) = \text{odd part of } (3 \cdot 0 + 1) = \text{odd part of } 1 = 1$.
\end{proof}

\begin{lemma}\label{lemma:R_one}\lean{CollatzMapBasics.R_one}\leanok
$R(1) = 1$.
\end{lemma}
\begin{proof}\leanok
By definition, $R(1) = \text{odd part of } (3 \cdot 1 + 1) = \text{odd part of } 4 = 1$.
\end{proof}

\begin{lemma}\label{lemma:reduced_collatz_step_pos}\lean{CollatzMapBasics.reduced_collatz_step_pos}\leanok
For all $n \in \mathbb{N}$, $R(n) \geq 1$.
\end{lemma}
\begin{proof}\leanok
Since $3n + 1 \geq 1$ for all $n$, the odd part of $3n+1$ is a positive divisor of a positive number, hence $R(n) \geq 1$.
\end{proof}

\begin{lemma}\label{lemma:R_pos}\lean{CollatzMapBasics.R_pos}\leanok\uses{lemma:reduced_collatz_step_pos}
For all $n, i \in \mathbb{N}$ with $n \geq 1$, we have $R^i(n) \geq 1$.
\end{lemma}
\begin{proof}\leanok
We proceed by induction on $i$. The base case $i = 0$ is immediate since $R^0(n) = n \geq 1$. For the inductive step, $R^{i+1}(n) = R^i(R(n))$. By Lemma~\ref{lemma:reduced_collatz_step_pos}, $R(n) \geq 1$, so the induction hypothesis gives $R^i(R(n)) \geq 1$.
\end{proof}

\begin{lemma}\label{lemma:reduced_collatz_step_odd}\lean{CollatzMapBasics.reduced_collatz_step_odd}\leanok
For all $n \in \mathbb{N}$, $R(n)$ is odd, i.e.\ $R(n) \equiv 1 \pmod{2}$.
\end{lemma}
\begin{proof}\leanok
By definition, $R(n)$ is the odd part of $3n+1$.
\end{proof}

\begin{lemma}\label{lemma:R_odd}\lean{CollatzMapBasics.R_odd}\leanok\uses{lemma:reduced_collatz_step_odd}
For all $n \in \mathbb{N}$ and $i \geq 1$, $R^i(n)$ is odd, i.e.\ $R^i(n) \equiv 1 \pmod{2}$.
\end{lemma}
\begin{proof}\leanok
We proceed by induction on $i$. The base case $i = 1$ follows directly from Lemma~\ref{lemma:reduced_collatz_step_odd}, since $R^1(n) = R(n)$ is odd. For the inductive step, $R^{i+1}(n) = R^i(R(n))$. If $i \geq 1$, the induction hypothesis gives that $R^i(R(n))$ is odd. If $i = 0$, then $R^{i+1}(n) = R^1(n) = R(n)$, which is odd by Lemma~\ref{lemma:reduced_collatz_step_odd}.
\end{proof}

\begin{lemma}\label{lemma:not_exists_R_with_m_div3}\lean{CollatzMapBasics.not_exists_R_with_m_div3}\leanok
If $m \equiv 0 \pmod{3}$, then there is no $n \in \mathbb{N}$ such that $R(n) = m$. In other words, multiples of $3$ are not in the image of $R$.
\end{lemma}
\begin{proof}\leanok
Suppose for contradiction that $R(n) = m$ for some $n$, with $3 \mid m$. Since $R(n)$ is the odd
part of $3n+1$, we have $R(n) \mid 3n+1$.
\end{proof}

\begin{lemma}\label{lemma:exists_R_with_m_not_div3}\lean{CollatzMapBasics.exists_R_with_m_not_div3}\leanok
Let $m > 0$ be odd with $m \not\equiv 0 \pmod{3}$. Then there exists an odd $n > 1$ such that $R(n) = m$.
\end{lemma}
\begin{proof}\leanok
If $m = 1$, take $n = 5$; one checks $R(5) = 1$ by computation.

If $m > 1$, we split on the residue of $m$ modulo $3$:
\begin{itemize}
\item If $m \equiv 1 \pmod{3}$: set $n = (4m - 1)/3$. Then $3 \mid (4m-1)$ and $3n + 1 = 4m = 2^2 \cdot m$. Since $m$ is odd, the odd part of $2^2 \cdot m$ is $m$, so $R(n) = m$. One verifies $n$ is odd and $n > 1$.
\item If $m \equiv 2 \pmod{3}$: set $n = (2m - 1)/3$. Then $3 \mid (2m-1)$ and $3n + 1 = 2m = 2^1 \cdot m$. Since $m$ is odd, the odd part of $2m$ is $m$, so $R(n) = m$. One verifies $n$ is odd and $n > 1$.
\end{itemize}
\end{proof}

\begin{lemma}\label{lemma:reduced_collatz_step_iff_mod4}\lean{CollatzMapBasics.reduced_collatz_step_iff_mod4}\leanok
For any odd $n > 0$, $R(n) > n$ if and only if $n \equiv 3 \pmod{4}$.
\end{lemma}
\begin{proof}\leanok
($\Rightarrow$) Suppose $n \equiv 1 \pmod{4}$, say $n = 4k + 1$. Then $3n + 1 = 12k + 4 = 2^2(3k + 1)$. The reduced step is $R(n) = \text{odd part of } (3k + 1) \le 3k + 1$. Since $3k + 1 < 4k + 1$ (for $k \ge 0$), we have $R(n) < n$, a contradiction. Thus we must have $n \equiv 3 \pmod{4}$.

($\Leftarrow$) Suppose $n \equiv 3 \pmod{4}$, say $n = 4k + 3$. Then $3n + 1 = 12k + 10 = 2(6k + 5)$. Since $6k + 5$ is odd, $R(n) = 6k + 5$. Since $k \ge 0$, we have $6k + 5 > 4k + 3$, so $R(n) > n$.
\end{proof}

\begin{lemma}\label{lemma:reduced_collatz_step_implies_collatz_iter}\lean{CollatzMapBasics.reduced_collatz_step_implies_collatz_iter}\leanok
If $n > 0$ is odd and $R(n) = m$, then there exists $i \ge 1$ such that $C^i(n) = m$.
\end{lemma}
\begin{proof}\leanok
By definition, $R(n)$ is the odd part of $3n+1$, so $3n+1 = 2^k \cdot m$ for some $k \in \mathbb{N}$. Since $n$ is odd, $3n+1$ is even, so $k \ge 1$.
One iteration of the Collatz map on $n$ gives $C^1(n) = 3n+1 = 2^k \cdot m$.
By Lemma~\ref{lemma:collatz_iter_two_pow_mul}, applying the Collatz map $k$ more times to $2^k \cdot m$ yields $C^k(2^k \cdot m) = m$.
Thus $C^{k+1}(n) = C^k(C(n)) = C^k(2^k \cdot m) = m$.
\end{proof}

\begin{lemma}\label{lemma:reduced_collatz_iter_implies_collatz_iter}\lean{CollatzMapBasics.reduced_collatz_iter_implies_collatz_iter}\leanok
If $n$ is odd and $R^k(n) = m$, then there exists $i \ge k$ such that $C^i(n) = m$.
\end{lemma}
\begin{proof}\leanok
We proceed by induction on $k$. The base case $k=0$ is trivial as $R^0(n) = n = C^0(n)$.
For the inductive step $k+1$, let $n_1 = R(n)$. By the induction hypothesis, since $R^k(n_1) = m$, there exists $i \ge k$ such that $C^i(n_1) = m$.
By Lemma~\ref{lemma:reduced_collatz_step_implies_collatz_iter}, there exists $j \ge 1$ such that $C^j(n) = n_1$.
Then $C^{i+j}(n) = C^i(C^j(n)) = C^i(n_1) = m$.
Since $i \ge k$ and $j \ge 1$, we have $i+j \ge k+1$.
\end{proof}

\begin{definition}\label{def:count_odds}\lean{CollatzMapBasics.count_odds}\leanok
For $d, n \in \mathbb{N}$, the number of odd steps in the first $d$ iterations of the standard Collatz map starting from $n$ is defined by
\begin{equation}
    \text{countOdds}(d, n) = \sum_{j=0}^{d-1} [C^j(n) \text{ is odd}]
\end{equation}
where $[P]$ is the Iverson bracket.
\end{definition}

\begin{lemma}\label{lemma:ordCompl_collatz_iter}\lean{CollatzMapBasics.ordCompl_collatz_iter}\leanok
Let $n \in \mathbb{N}$ be odd. For any $d \in \mathbb{N}$, the odd part of $C^d(n)$ is equal to $R^k(n)$, where $k = \text{countOdds}(d, n)$. That is,
\begin{equation}
    \text{odd part of } C^d(n) = R^k(n).
\end{equation}
\end{lemma}
\begin{proof}\leanok
We proceed by induction on $d$. The base case $d=0$ is trivial as $C^0(n) = n$ is odd, its odd part is $n$, and $\text{count\_odds}(0, n) = 0$, with $R^0(n) = n$.

For the inductive step $d+1$, let $C^d(n) = m$. By induction, the odd part of $m$ is $R^{\text{countOdds}(d, n)}(n)$.
If $m$ is even, then $C^{d+1}(n) = m/2$. The odd part of $m/2$ is the same as the odd part of $m$.
Since $m$ is even, $\text{countOdds}(d+1, n) = \text{countOdds}(d, n)$, and the equality holds.
If $m$ is odd, then $C^{d+1}(n) = 3m+1$. Since $m$ is odd, it equals its own odd part, $m = R^{\text{countOdds}(d, n)}(n)$.
The odd part of $3m+1$ is $R(m) = R(R^{\text{countOdds}(d, n)}(n)) = R^{\text{countOdds}(d, n) + 1}(n)$.
Since $m$ is odd, $\text{countOdds}(d+1, n) = \text{countOdds}(d, n) + 1$, confirming the result.
\end{proof}

\begin{lemma}\label{lemma:collatz_cycle_iff_reduced_collatz_cycle}\lean{CollatzMapBasics.collatz_cycle_iff_reduced_collatz_cycle}\leanok
Let $n \in \mathbb{N}$ be odd. There exists a cycle in the standard Collatz map starting from $n$ if and only if there exists a cycle in the reduced Collatz map starting from $n$.
\end{lemma}
\begin{proof}\leanok
($\Rightarrow$) Suppose there exists $k > 0$ such that $C^k(n) = n$. By Lemma~\ref{lemma:ordCompl_collatz_iter}, the odd part of $C^k(n)$ is $R^j(n)$ where $j = \text{countOdds}(k, n)$. Since $n$ is odd, its odd part is $n$, so $R^j(n) = n$. Since $n$ is odd and $k > 0$, at least one standard Collatz step must have occurred, and since $n \ge 1$, we must have $j > 0$. Thus $n$ is part of a reduced cycle.

($\Leftarrow$) Suppose there exists $j > 0$ such that $R^j(n) = n$. By Lemma~\ref{lemma:reduced_collatz_iter_implies_collatz_iter}, there exists $k \ge j$ such that $C^k(n) = n$. Since $j > 0$, we have $k > 0$, so $n$ is part of a standard cycle.
\end{proof}

\begin{lemma}\label{lemma:collatz_bounded_iff_reduced_collatz_bounded}\lean{CollatzMapBasics.collatz_bounded_iff_reduced_collatz_bounded}\leanok
Let $n \in \mathbb{N}$ be odd. The orbit of $n$ under the standard Collatz map is bounded if and only if its orbit under the reduced Collatz map is bounded.
\end{lemma}
\begin{proof}\leanok
($\Rightarrow$) Suppose there exists $B$ such that $C^k(n) \le B$ for all $k \ge 0$. By Lemma~\ref{lemma:reduced_collatz_iter_implies_collatz_iter}, every reduced iterate $R^j(n)$ is equal to some $C^k(n)$, so $R^j(n) \le B$ for all $j \ge 0$.

($\Leftarrow$) Suppose there exists $B$ such that $R^j(n) \le B$ for all $j \ge 0$. We claim that $C^k(n) \le \max(n, 3B+1)$ for all $k \ge 0$. We proceed by induction on $k$. The base case $k=0$ is $C^0(n) = n$, which is clearly bounded.
For the inductive step, let $C^k(n) = m$.
If $m$ is even, $C^{k+1}(n) = m/2 \le m$, which is bounded by the induction hypothesis.
If $m$ is odd, then by Lemma~\ref{lemma:ordCompl_collatz_iter}, $m$ is the odd part of $C^k(n)$, which is $R^j(n)$ for $j = \text{countOdds}(k, n)$. Thus $m \le B$. Then $C^{k+1}(n) = 3m+1 \le 3B+1$, which is within the bound.
\end{proof}
